\section{Biology Access, Verified}
\sectionkicker{VERIFIED ACCESS}
\begin{wideflow}
{\Large\bfseries 生物学への扉、検証つきで開く}\par
\smallskip
{\color{SurveyMuted}9月17日のLSVPベータを、利用枠と監視の仕組みとともに示す。}\par
\end{wideflow}

9月17日、Anthropicは生命科学向けの検証つき利用制度「Life Sciences Verification Program（LSVP）」をベータで開いた\cite{lsvp}。汎用のFableでは塞がれていた創薬・研究生物学・臨床開発・製造の仕事に、Mythos・Opus・Sonnet級を生物学向けに緩めた安全装置で触れるようにする内容で、まず団体・機関向け、個人のPro・Maxは後になるとされる。申請者は研究資格・安全基準・倫理審査の確認を通る。標準利用は団体全体に年更新で、基礎研究から開発・供給網・製造、臨床開発・品質・薬事、投資判断までを覆う。高リスク利用は一つの研究計画に限る追加枠で半年更新、ウイルスベクター一家の免疫認識の特徴づけのような申込みが例示される。Mythosの高リスク枠は小規模に絞り、米政府との調整を待つ追加審査つきとされる。

運用の要は、実時間の遮断に代わる事後監視である\cite{lsvp}。目印のついた活動の記録を30日保ち、厳密に隔てて学習には使わず、生命科学の研究側からは触れさせない。利用は申告の用途に結びつけ、範囲外の振る舞いは団体の管理者に事前の取り決めの時間枠で知らせるとされる。狙われる形として、侵入・乗っ取り、内部の不正・強要、群れや長期のエージェント悪用が挙げられる。提供面は一次のコンソール（API）とClaude Enterprise・Teamで、外部基盤には載らず、ベータはBAA対象の団体向けではない。最初の週に数百団体を見込むとされるが、加入や保持の見立てはベンダー側の計画として受け止める。二重用途の危うさの語りはベンダー側の文脈である。

\begin{claimboundary}[制度と公開日の境界]
Mythos 5.1の名称はこの制度頁の適用範囲に現れるが、Mythos 5.1の公開日の根拠にはしない。加入や保持の見立て、二重用途の語りはベンダー側の計画・文脈として扱う。
\end{claimboundary}

\begin{communitynote}[X / Community observation --- context only]
公式発信のLSVP告知や研究開発指標の投稿は、発表の存在を示す合図として扱う\cite{w38x-lsvp-off,w38x-rdi-off}。制度の仕組みや安全性の事実は発表資料そのものに拠り、公開投稿をもって技術の裏づけとはしない。
\end{communitynote}
